\documentclass[12pt,a4paper]{article}

\input{/home/tabaka/Documents/GitHub/Defs/defs.tex}

\setcounter{zadaniaRozdzial}{1}

\begin{document}
		
	\section{Liczby rzeczywiste}
	
	\subsection{Liczby naturalne}\:
	
	\ZbiorZadanieTwoXThree{Wyznacz dzielniki liczby:}{54}{21}{50}{125}{42}{40}
	
	\ZbiorZadanieTwoXTwo{Zapisz podaną liczbę w postaci $2k+r$, gdzie $r<k$:}{51}{20}{31}{17}
	
	\ZbiorZadanieTwoXTwo{Zapisz podaną liczbę w postaci $7k+r$, gdzie $r<k$:}{15}{124}{68}{84}
	
	\ZbiorZadanieTwoXTwo{Zapisz podaną liczbę w postaci $9k+r$, gdzie $r<k$:}{19}{120}{81}{54}
	
	\Wzor{Liczbę naturalną, która ma dokładnie dwa dzielniki (1 i samą siebie), nazywamy \textbf{liczbą pierwszą}. Każdą liczbę, która ma więcej niż te dwa dzielniki, nazywamy \textbf{liczbą złożoną}.}
	
	\ZbiorZadanieTwoXThree{Rozłóż liczby na czynniki pierwsze:}{100}{64}{150}{770}{294}{123}
	
	\ZbiorZadanieTwoXThree{Wyznacz NWD i NWW liczb:}{72 i 48}{18 i 30}{15 i 50}{144 i 192}{84 i 105}{196 i 420}
	
	\newpage
	
	\ZbiorZadanie{Udowodnij, że:
		\begin{enumerate}[a)]
			\item suma trzech kolejnych liczb naturalnych jest podzielna przez 3.
			\item suma trzech kolejnych liczb parzystych jest podzielna przez 6.
			\item suma dwóch kolejnych liczb nieparzystych jest podzielna przez 4.
	\end{enumerate}}
	
	\ZbiorZadanie{Udowodnij, że:
		\begin{enumerate}[a)]
			\item iloczyn dwóch kolejnych liczb naturalnych jest podzielny przez 2.
			\item iloczyn trzech kolejnych liczb naturalnych jest podzielny przez 6.
			\item iloczyn trzech kolejnych liczb parzystych jest podzielny przez 48.
	\end{enumerate}}
	
	\subsection{Liczby całkowite}\:
	
	\ZbiorZadanieTwoXTwo{Oblicz:}{$4-(6-(12-10))$}{$-9+((14-6)-(7-13))$}{$7-(4+(-12 -(-4)))$}{$13(5-19)-3(9-13)$}
	
	\ZbiorZadanieTwoXTwo{Oblicz:}
	{$18-(12-(7+5))\cdot2$}
	{$-8+((15-9)+(6-2))\cdot3$}
	{$4((13-17)-6(19-11)):2$}
	{$(42:(-6-8))((-5-17):(5-7))$}
	
	\subsection{Liczby wymierne}\:
	
	\ZbiorZadanieTwoXThree{Oblicz:}{$1\frac{2}{5}-2\frac{1}{6}$}{$3\frac{3}{8}+2\frac{5}{6}$}{$\frac{7}{6}+(-\frac{5}{4}-\frac{4}{3})$}{1$\frac{5}{8}-(-2\frac{2}{3}+\frac{1}{4})$}{$\frac{1}{2}-\frac{1}{4}-(\frac{1}{4}-\frac{1}{5})$}{$-1\frac{9}{16}-\frac{11}{12}$}
	
	\ZbiorZadanieTwoXTwo{Oblicz:}{$-1\frac{2}{25}\cdot (-2\frac{1}{12})$}{4$\frac{1}{7} - 4\frac{1}{2}:(-5\frac{1}{4})$}{$-2\frac{1}{9}\cdot 1\frac{16}{19}:\frac{7}{5}$}{$2\frac{1}{3}:2\frac{4}{5}\cdot\frac{18}{77}\cdot2\frac{17}{30}$}
	
	\ZbiorZadanieTwoXTwo{Oblicz:}{$\LARGE\frac{2\frac{1}{6}}{1\frac{4}{9}}$}{$\frac{\frac{1}{2}+\frac{1}{3}}{\frac{1}{2}}$}{$\frac{2-\frac{1}{3}}{2-\frac{1}{3}}$}{$\frac{2\frac{1}{4}-3\frac{1}{2}}{\frac{1}{8}+\frac{1}{2}}$}
	
	\ZbiorZadanieTwoXThree{Oblicz:}
	{$2,45+3,7$}
	{$8,6+1,275$}
	{$4,08+5,92$}
	{$9,5-3,75$}
	{$12,04-6,8$}
	{$7,3-2,456$}
	
	\ZbiorZadanieTwoXThree{Oblicz:}
	{$2,4\cdot3,5$}
	{$0,75\cdot4,8$}
	{$1,25\cdot0,64$}
	{$7,2\cdot1,5$}
	{$8,4:0,7$}
	{$12,6:1,4$}
	
	\ZbiorZadanieTwoXTwo{Zapisz w postaci ułamka dziesiętnego:}{$\frac{1}{3}$}{$2\frac{2}{7}$}{$\frac{11}{6}$}{$\frac{7}{90}$}
	
	\ZbiorZadanieTwoXThree{Zapisz w postaci ułamka zwykłego:}
	{$0,(3)$}
	{$0,1(6)$}
	{$1,(25)$}
	{$2,3(4)$}
	{$0,12(7)$}
	{$0,(9)$}
	
	\subsection{Działania na pierwiastkach}\:
	
	\ZbiorZadanieTwoXThree{Oblicz:}
	{$\sqrt{12}+\sqrt{27}$}
	{$3\sqrt{5}-\sqrt{20}$}
	{$\sqrt{6}\cdot\sqrt{24}$}
	{$\frac{\sqrt{75}}{\sqrt{3}}$}
	{$2\sqrt{18}+\sqrt{50}$}
	{$\sqrt{48}-\sqrt{12}$}
	
	\ZbiorZadanieTwoXThree{Oblicz:}
	{$2\sqrt{3}+4\sqrt{3}-\sqrt{27}$}
	{$5\sqrt{2}-\sqrt{8}+2\sqrt{18}$}
	{$(\sqrt{3}+2)(\sqrt{3}-2)$}
	{$\sqrt{0,25}+\sqrt{1,44}+\sqrt{6,25}$}
	{$\frac{6\sqrt{15}}{3\sqrt{5}}$}
	{$\sqrt{8}\left(\sqrt{2}+\sqrt{18}\right)$}
	
	\ZbiorZadanie{Podaj wszystkie liczby naturalne leżące między $\sqrt{10}$ i $\sqrt{140}$.}
	
	\ZbiorZadanieTwoXThree{Oblicz:}
	{$\sqrt[3]{27}$}
	{$\sqrt[3]{-64}$}
	{$\sqrt[4]{81}$}
	{$\sqrt[5]{32}$}
	{$\sqrt[3]{8}+\sqrt[3]{125}$}
	{$\sqrt[3]{216}-\sqrt[3]{27}$}
	
	\ZbiorZadanieTwoXThree{Wyłącz czynnik przed znak pierwiastka:}
	{$\sqrt[3]{54}$}
	{$\sqrt[3]{128}$}
	{$\sqrt[4]{48}$}
	{$\sqrt[4]{80}$}
	{$\sqrt[5]{96}$}
	{$\sqrt[3]{250}$}
	
	\subsection{Działania na potęgach}\:
	
	\ZbiorZadanieTwoXThree{Oblicz:}
	{$2^3+2^4$}
	{$5^2-3^2$}
	{$2^3\cdot2^2$}
	{$3^5:3^2$}
	{$(2^3)^2$}
	{$5^0+2^3$}
	
	\ZbiorZadanieTwoXThree{Oblicz:}
	{$2^4\cdot2^3$}
	{$5^6:5^2$}
	{$3^2\cdot3^4:3^3$}
	{$7^5:7^3\cdot7^2$}
	{$(2^3)^4$}
	{$(5^2)^3:5^4$}
	
	\ZbiorZadanieTwoXThree{Oblicz:}
	{$(-2)^4$}
	{$(-3)^3$}
	{$(-2)^5+2^5$}
	{$\left(\frac{1}{2}\right)^3$}
	{$\left(-\frac{2}{3}\right)^2$}
	{$\left(\frac{3}{4}\right)^2:\frac{3}{4}$}
	
	\ZbiorZadanieTwoXThree{Oblicz:}
	{$2^3\cdot2^4:2^5$}
	{$3^5:3^2+2^3$}
	{$(2^2)^3:2^4$}
	{$5^3-2\cdot5^2$}
	{$2^4\cdot3^2:2^2$}
	{$(3^2+3^3):3^2$}
	
	\ZbiorZadanieTwoXThree{Oblicz:}
	{$\left(\frac{2}{3}\right)^2$}
	{$\left(\frac{3}{4}\right)^3$}
	{$\left(\frac{5}{2}\right)^2$}
	{$\left(-\frac{2}{5}\right)^3$}
	{$\left(\frac{4}{3}\right)^2$}
	{$\left(-\frac{3}{2}\right)^2$}
	
	\ZbiorZadanieTwoXThree{Oblicz:}
	{$\frac{2^4}{2^2}$}
	{$\frac{3^5}{3^2}$}
	{$\frac{5^3}{5}$}
	{$\frac{2^3\cdot2^2}{2^4}$}
	{$\frac{3^4\cdot3^2}{3^3}$}
	{$\frac{5^4}{5^2\cdot5}$}
	
	\ZbiorZadanieTwoXThree{Oblicz:}
	{$\frac{6^3}{8\cdot9^2}$}
	{$\frac{12^2}{2^3\cdot3^2}$}
	{$\frac{18^2}{4\cdot9^2}$}
	{$\frac{20^2}{2^3\cdot5^2}$}
	{$\frac{15^3}{3^2\cdot5^2}$}
	{$\frac{24^2}{2^4\cdot3^2}$}
	
	\subsection{Potęga o wykładniku wymiernym}\:
	
	\ZbiorZadanieTwoXThree{Oblicz:}
	{$2^{-3}$}
	{$5^{-2}$}
	{$\left(\frac{2}{3}\right)^{-2}$}
	{$(-3)^{-3}$}
	{$\frac{4^{-2}}{2^{-3}}$}
	{$\left(\frac{1}{5}\right)^{-2}$}
	
	\ZbiorZadanieTwoXThree{Oblicz:}
	{$8^{\frac{1}{3}}$}
	{$16^{\frac{1}{2}}$}
	{$27^{\frac{2}{3}}$}
	{$32^{\frac{2}{5}}$}
	{$81^{\frac{3}{4}}$}
	{$\left(\frac{1}{16}\right)^{\frac{3}{4}}$}
	
	\ZbiorZadanieTwoXThree{Oblicz:}
	{$\sqrt[3]{5\sqrt{5}}$}
	{$\sqrt{2\sqrt{2}}$}
	{$\sqrt[3]{4\sqrt{2}}$}
	{$\sqrt[4]{8\sqrt{2}}$}
	{$\sqrt{3\sqrt{3}}$}
	{$\sqrt[3]{9\sqrt{3}}$}
	
\end{document}